\section{Orientation}
\label{part:orientation}

\subsection{The statement}

Fix a compiled artifact $P$. Partition its inputs into public inputs $p$ and
secret inputs $s$. Fix an observer and let $\Obs_\Theta$ be the projection of an
execution trace onto what that observer can actually see under target profile
$\Theta$ --- branch directions, memory addresses, operation latencies and
counts, allocation sizes, transfer lengths, values arriving at public sinks, and
the identity of the host or audience a message is delivered to.

Some dependence on the secret is intended: a service that returns nothing
derived from its inputs is useless. So the property is not plain
noninterference but \emph{release-relative} noninterference. Let $\Rel$ be the
releases authorized for that observer. Then:
%
\begin{equation}
\label{eq:ni}
\forall p,\, s_0,\, s_1 :\quad
  \Rel(p,s_0) = \Rel(p,s_1)
  \;\Longrightarrow\;
  \Obs_\Theta\bigl(\Trace(P,p,s_0)\bigr) = \Obs_\Theta\bigl(\Trace(P,p,s_1)\bigr).
\end{equation}
%
Read it as a challenge from an adversary. They choose the public inputs once and
two secrets freely. They only have to respect one constraint: whatever they are
\emph{entitled} to learn must come out the same both times. If they can then
make any observable differ, they have learned something they were not entitled
to, and $P$ leaks.

\begin{remark}
Equation~\eqref{eq:ni} quantifies over \emph{pairs} of executions. That single
fact drives almost every design decision in the rest of this document, and
Part~\ref{part:background} is devoted to its consequences.

Two warnings about the form above, both discharged there. It is a statement of
\emph{intent}, not an encoding: installing its antecedent as an initial
constraint on the pair is a known unsoundness, and \cref{sec:ledger} gives the
prefix-causal construction that replaces it. And $P$ here is informal shorthand;
from \cref{step:statement} onward the object under analysis is the frozen
artifact $\artifact$ together with its pipeline $\pipeline$, because a verdict is
only ever a statement about specific bytes.
\end{remark}

\subsection{Why functional correctness is not enough}

A verified compiler such as CompCert proves that the target program computes the
same \emph{function} as the source. That is a statement about one execution at a
time, and confidentiality is not a property of one execution at a time.

Concretely, every one of the following transformations preserves the computed
function and destroys \eqref{eq:ni}:

\begin{itemize}[nosep,leftmargin=1.4em]
\item replacing a constant-time conditional move with a branch, so the observer
      learns the condition from control flow;
\item strength-reducing a division into a multiply-high, changing which
      operations are executed and how long they take;
\item eliminating a store that was scrubbing a secret, because nothing reads it
      afterwards;
\item promoting a secret buffer into registers that are later spilled;
\item choosing an allocation size that depends on a secret count.
\end{itemize}

So a confidentiality checker cannot be a corollary of a correctness checker. It
needs its own property, its own observation model, and --- because the property
is about pairs --- its own machinery.

\subsection{Four outcomes, and why three would be dishonest}

The checker reports one of four values for each question it is asked:

\begin{center}
\begin{tabular}{@{}l p{0.44\linewidth} p{0.27\linewidth}@{}}
\toprule
Outcome & Meaning & Obligations \\
\midrule
\Verified    & no residual forbidden dependence at the claimed boundary
             & must be none \\
\Unsafe      & a forbidden dependence, witnessed by a replayable counterexample
             & must be none \\
\Unknown     & required semantics or a contract is absent
             & must name what is missing \\
\Conditional & acceptable only if named external assumptions hold
             & must name every assumption \\
\bottomrule
\end{tabular}
\end{center}

The two middle-ground values carry the weight. A checker that can only say
\emph{safe} or \emph{unsafe} must guess when it meets a function pointer, a
recursive call, an opaque external callee, an uninitialized read, or a solver
timeout --- and a guess in the \Verified{} direction is exactly the failure that
makes such a tool worthless. \Unknown{} is what lets the analysis fail closed.
\Conditional{} is what lets it be honest about facts that live outside the model
altogether: real helper latency, backend trace preservation, hardware isolation.

\begin{remark}[The invariant that makes the ledger work]
\Verified{} and \Unsafe{} must carry \emph{no} outstanding obligations, while
\Unknown{} and \Conditional{} must name theirs. That is precisely the discipline
a claim ledger encodes: an outcome is only allowed to be terminal when the
ledger for it is empty.
\end{remark}

\subsection{Evidence levels}

The system separates \emph{where} a fact is established from \emph{how much} it
is worth. Five levels are used throughout, and every fixture in the harness
names the ones it relies on.

\begin{center}
\begin{tabular}{@{}l p{0.80\linewidth}@{}}
\toprule
Level & What it supplies \\
\midrule
\Lv{0} & declarations: secrets, observer, release policies, helper contracts,
         hosts, target facts \\
\Lv{1} & propagation: dependence and effects --- branches, addresses,
         latencies, sinks \\
\Lv{2} & a relational witness: two lanes coupled at entry, different secrets,
         different observations \\
\Lv{3} & attribution: relating source, imported IR, and target artifact to a
         compiler boundary \\
\Lv{4} & facts outside the modeled semantics: real latency, compressors,
         hardware isolation, certificates \\
\bottomrule
\end{tabular}
\end{center}

\Lv{0} is a \emph{contract, not a proof of the contract}. If no latency summary
for a helper is available, the honest result is \Unknown{}; selecting a named
operand-dependent test profile lets \Lv{1} classify the call as \Unsafe{} under
that profile, and \Lv{4} is still required before claiming a deployed helper
matches it.

\subsection{How to read the ledger}
\label{sec:ledgerguide}

Each step in Part~\ref{part:walkthrough} carries a panel in the right margin
listing the obligations \emph{still open} after that step. Rows are coloured by
what the step did:

\begin{center}
\begin{tabular}{@{}ll@{}}
\toprule
\textsc{new} (green) & the step opened a new obligation \\
\textsc{reduced} (amber) & the step rewrote an obligation into a different one \\
\textsc{used} (amber) & the step relied on it; the body is unchanged \\
\textsc{discharged} (red) & the step closed it; struck out once, then dropped \\
carried (grey) & untouched by this step, still open \\
\bottomrule
\end{tabular}
\end{center}

Two reading conventions. First, an obligation quantified over an index --- per
coalition, per site, per entry --- is \emph{one} row carrying a subscript, not
one row per instance; otherwise the panel would be unreadable by Step~3.
Second, the bodies are deliberately terse because a struck-out body cannot break
across lines. The abbreviations are:

\begin{center}
\begin{tabular}{@{}ll@{}}
\toprule
$\Sec_\coal(P)$ & $P$ satisfies \eqref{eq:ni} for coalition $\coal$ \\
$\Bind(\artifact,\pipeline)$ & every result is bound to this exact artifact and pipeline \\
$\Total(\manifest)$ & the policy manifest is total over reachable sites \\
$\NFConf(\artifact,\pipeline)$ & the artifact lies inside the analyzable fragment \\
$\Adm \neq \emptyset$ & at least one execution satisfies the declared preconditions \\
$\PairDom_\coal \neq \emptyset$ & at least one admitted low-equal \emph{pair} exists \\
$\DiagOK$ & the diagnostic layer assigned a disposition to every site \\
$\ProdSafe_\coal$ & the exact two-lane product cannot reach a bad state \\
$\ReplayCov_\coal(w)$ & witness $w$ replays under the exact semantics \\
$\PairedRef_\coal$ & the backend preserves the observable trace \\
\bottomrule
\end{tabular}
\end{center}

The arc opens with \eqref{eq:ni} itself and ends with an empty ledger. If you
only read the margins, you will still see which obligation each step traded
away and what it bought.
